\documentclass[a4paper,11pt]{ltjsarticle}
\usepackage{luatexja}
\usepackage{amsmath,amssymb,amsthm}
\usepackage{hyperref}
\usepackage[margin=25mm]{geometry}

\newtheorem{theorem}{定理}[section]
\newtheorem{proposition}[theorem]{命題}
\newtheorem{lemma}[theorem]{補題}
\newtheorem{corollary}[theorem]{系}
\newtheorem{definition}[theorem]{定義}
\newtheorem{remark}[theorem]{注意}
\newtheorem{observation}[theorem]{観察}

\providecommand{\tightlist}{\setlength{\itemsep}{0pt}\setlength{\parskip}{0pt}}

\begin{document}

\section{標準模型の19個の任意パラメータの分類と構造分析}\label{ux6a19ux6e96ux6a21ux578bux306e19ux500bux306eux4efbux610fux30d1ux30e9ux30e1ux30fcux30bfux306eux5206ux985eux3068ux69cbux9020ux5206ux6790}

\subsection{── 万物の理論の構造要件に基づく評価
──}\label{ux4e07ux7269ux306eux7406ux8ad6ux306eux69cbux9020ux8981ux4ef6ux306bux57faux3065ux304fux8a55ux4fa1}

\textbf{著者:} 木原 範昭 (WF System Co., Ltd.)

\textbf{日付:} 2026年4月

\textbf{DOI:}
\href{https://doi.org/10.5281/zenodo.19604965}{10.5281/zenodo.19604965}

\textbf{前提論文:}
\href{https://doi.org/10.5281/zenodo.19601592}{万物の理論が満たすべき構造要件の定式化}

\begin{center}\rule{0.5\linewidth}{0.5pt}\end{center}

\subsection{0. 本稿の目的}\label{ux672cux7a3fux306eux76eeux7684}

本稿は、素粒子物理学の標準模型に含まれる19個の任意パラメータ（自由パラメータ）を、前稿で定式化した構造要件の観点から分類・分析する。

具体的には、以下の問いに答える。

\begin{enumerate}
\def\labelenumi{\arabic{enumi}.}
\tightlist
\item
  各パラメータは、構造的に導出可能か、それとも真に任意か。
\item
  万物の理論が内部的に生成すべきパラメータはどれか（= TOEのKPI）。
\item
  パラメータ間の依存関係はどうなっているか。
\end{enumerate}

\begin{center}\rule{0.5\linewidth}{0.5pt}\end{center}

\subsection{1.
パラメータ一覧}\label{ux30d1ux30e9ux30e1ux30fcux30bfux4e00ux89a7}

現代物理学の基本理論（標準模型 +
一般相対性理論）に関わるパラメータの全体像を示す。標準模型の慣例では「任意パラメータ」は19個とされるが、万物の理論の評価にあたっては、「基本定数」として固定されているものも含め、すべてのパラメータを検証対象とする。

各パラメータの値は、現在の実験で確認されている低エネルギー領域（概ね電子ボルト〜TeVスケール、すなわち原子・素粒子実験で到達可能な範囲）における値である。

\subsubsection{1.0
基本物理定数（現在「定数」とされているもの）}\label{ux57faux672cux7269ux7406ux5b9aux6570ux73feux5728ux5b9aux6570ux3068ux3055ux308cux3066ux3044ux308bux3082ux306e}

以下は現在の物理学では定数として扱われ、自然単位系（\(c = \hbar = k_B = 1\)）では消去されるため「任意パラメータ」にはカウントされない。しかし、大統一理論や量子重力理論の候補の中には、これらが宇宙の歴史の中で変化する、あるいはより基本的な量から導出されるとする仮説も存在する。万物の理論の検証対象として一覧に含める。

これらの定数は、関与する次元の構成によって以下のグループに分類される。この分類は、各定数が物理学のどの層に属するかを明確にする。

\paragraph{\texorpdfstring{グループ
I：純粋時空（\(L, T\)）}{グループ I：純粋時空（L, T）}}\label{ux30b0ux30ebux30fcux30d7-iux7d14ux7c8bux6642ux7a7al-t}

{\def\LTcaptype{none} % do not increment counter
\begin{longtable}[]{@{}
  >{\raggedright\arraybackslash}p{(\linewidth - 10\tabcolsep) * \real{0.0833}}
  >{\raggedright\arraybackslash}p{(\linewidth - 10\tabcolsep) * \real{0.3056}}
  >{\raggedright\arraybackslash}p{(\linewidth - 10\tabcolsep) * \real{0.1667}}
  >{\raggedright\arraybackslash}p{(\linewidth - 10\tabcolsep) * \real{0.1111}}
  >{\raggedright\arraybackslash}p{(\linewidth - 10\tabcolsep) * \real{0.1667}}
  >{\raggedright\arraybackslash}p{(\linewidth - 10\tabcolsep) * \real{0.1667}}@{}}
\toprule\noalign{}
\begin{minipage}[b]{\linewidth}\raggedright
\#
\end{minipage} & \begin{minipage}[b]{\linewidth}\raggedright
パラメータ
\end{minipage} & \begin{minipage}[b]{\linewidth}\raggedright
記号
\end{minipage} & \begin{minipage}[b]{\linewidth}\raggedright
値
\end{minipage} & \begin{minipage}[b]{\linewidth}\raggedright
次元
\end{minipage} & \begin{minipage}[b]{\linewidth}\raggedright
変換
\end{minipage} \\
\midrule\noalign{}
\endhead
\bottomrule\noalign{}
\endlastfoot
P1 & 光速度 & c & 299,792,458 m/s（定義値） & \([L T^{-1}]\) & 空間
\(L\) ↔ 時間 \(T\)：\(L = cT\) \\
\end{longtable}
}

\begin{itemize}
\tightlist
\item
  質量も電荷も温度も関与しない。\textbf{純粋な幾何学}の定数。
\item
  特殊相対性理論の基盤であり、時空を統合する唯一の定数。
\item
  TOEでの問い：時空の幾何学的構造から導出されるか？
  変光速理論（VSL）では可変とする仮説もある。
\end{itemize}

\paragraph{\texorpdfstring{グループ
II：時空＋質量（\(M, L, T\)）}{グループ II：時空＋質量（M, L, T）}}\label{ux30b0ux30ebux30fcux30d7-iiux6642ux7a7aux8ceaux91cfm-l-t}

{\def\LTcaptype{none} % do not increment counter
\begin{longtable}[]{@{}
  >{\raggedright\arraybackslash}p{(\linewidth - 10\tabcolsep) * \real{0.0833}}
  >{\raggedright\arraybackslash}p{(\linewidth - 10\tabcolsep) * \real{0.3056}}
  >{\raggedright\arraybackslash}p{(\linewidth - 10\tabcolsep) * \real{0.1667}}
  >{\raggedright\arraybackslash}p{(\linewidth - 10\tabcolsep) * \real{0.1111}}
  >{\raggedright\arraybackslash}p{(\linewidth - 10\tabcolsep) * \real{0.1667}}
  >{\raggedright\arraybackslash}p{(\linewidth - 10\tabcolsep) * \real{0.1667}}@{}}
\toprule\noalign{}
\begin{minipage}[b]{\linewidth}\raggedright
\#
\end{minipage} & \begin{minipage}[b]{\linewidth}\raggedright
パラメータ
\end{minipage} & \begin{minipage}[b]{\linewidth}\raggedright
記号
\end{minipage} & \begin{minipage}[b]{\linewidth}\raggedright
値
\end{minipage} & \begin{minipage}[b]{\linewidth}\raggedright
次元
\end{minipage} & \begin{minipage}[b]{\linewidth}\raggedright
変換
\end{minipage} \\
\midrule\noalign{}
\endhead
\bottomrule\noalign{}
\endlastfoot
P2 & プランク定数 & ℏ & 1.055 × 10⁻³⁴ J·s（定義値） &
\([M L^2 T^{-1}]\)（\(M^{+1}\)） & エネルギー \(E\) ↔ 周波数
\(\nu\)：\(E = \hbar\omega\) \\
P3 & 万有引力定数 & G & 6.674 × 10⁻¹¹ m³/(kg·s²) &
\([L^3 M^{-1} T^{-2}]\)（\(M^{-1}\)） & 質量 \(M\) ↔ 時空の曲率
\(L^{-2}\)：\(R \sim GM/(c^2 r)\) \\
\end{longtable}
}

\begin{itemize}
\tightlist
\item
  同じ \(M, L, T\) の3次元で構成されるが、\textbf{質量 \(M\)
  の符号が逆}。

  \begin{itemize}
  \tightlist
  \item
    ℏ：\(M^{+1}\) --- 質量が増えるとエネルギー（振動数）が増える
  \item
    G：\(M^{-1}\) ---
    質量が増えると重力定数の寄与は逆に小さくなる（曲率半径が大きくなる）
  \end{itemize}
\item
  この\textbf{質量に対する逆の振る舞い}が、量子力学と一般相対性理論が統合困難な根本原因を次元解析の段階で示唆している。
\item
  ℏ は量子力学（離散化）の基盤、G は一般相対性理論（幾何学化）の基盤。
\item
  G
  は次元変換定数でありながら重力の結合定数（力の強さ）でもある\textbf{二重性}を持つ。
\end{itemize}

\paragraph{\texorpdfstring{グループ
III：時空＋質量＋温度（\(M, L, T, \Theta\)）}{グループ III：時空＋質量＋温度（M, L, T, \textbackslash Theta）}}\label{ux30b0ux30ebux30fcux30d7-iiiux6642ux7a7aux8ceaux91cfux6e29ux5ea6m-l-t-theta}

{\def\LTcaptype{none} % do not increment counter
\begin{longtable}[]{@{}
  >{\raggedright\arraybackslash}p{(\linewidth - 10\tabcolsep) * \real{0.0833}}
  >{\raggedright\arraybackslash}p{(\linewidth - 10\tabcolsep) * \real{0.3056}}
  >{\raggedright\arraybackslash}p{(\linewidth - 10\tabcolsep) * \real{0.1667}}
  >{\raggedright\arraybackslash}p{(\linewidth - 10\tabcolsep) * \real{0.1111}}
  >{\raggedright\arraybackslash}p{(\linewidth - 10\tabcolsep) * \real{0.1667}}
  >{\raggedright\arraybackslash}p{(\linewidth - 10\tabcolsep) * \real{0.1667}}@{}}
\toprule\noalign{}
\begin{minipage}[b]{\linewidth}\raggedright
\#
\end{minipage} & \begin{minipage}[b]{\linewidth}\raggedright
パラメータ
\end{minipage} & \begin{minipage}[b]{\linewidth}\raggedright
記号
\end{minipage} & \begin{minipage}[b]{\linewidth}\raggedright
値
\end{minipage} & \begin{minipage}[b]{\linewidth}\raggedright
次元
\end{minipage} & \begin{minipage}[b]{\linewidth}\raggedright
変換
\end{minipage} \\
\midrule\noalign{}
\endhead
\bottomrule\noalign{}
\endlastfoot
P4 & ボルツマン定数 & k\_B & 1.381 × 10⁻²³ J/K（定義値） &
\([M L^2 T^{-2} \Theta^{-1}]\) & 温度 \(\Theta\) ↔ エネルギー
\(E\)：\(E = k_B T\) \\
\end{longtable}
}

\begin{itemize}
\tightlist
\item
  グループ II（\(M, L, T\)）に温度 \(\Theta\) の次元を追加する。
\item
  統計力学（多体系・マクロ）とミクロ物理を結ぶ定数。
\item
  情報のエントロピーと熱力学的エントロピーの橋渡し。情報理論的TOEでは
  \(k_B = 1\) が自然であり、独立パラメータではない可能性がある。
\end{itemize}

\paragraph{\texorpdfstring{グループ
IV：電荷のみ（\(Q\)）}{グループ IV：電荷のみ（Q）}}\label{ux30b0ux30ebux30fcux30d7-ivux96fbux8377ux306eux307fq}

{\def\LTcaptype{none} % do not increment counter
\begin{longtable}[]{@{}
  >{\raggedright\arraybackslash}p{(\linewidth - 10\tabcolsep) * \real{0.0833}}
  >{\raggedright\arraybackslash}p{(\linewidth - 10\tabcolsep) * \real{0.3056}}
  >{\raggedright\arraybackslash}p{(\linewidth - 10\tabcolsep) * \real{0.1667}}
  >{\raggedright\arraybackslash}p{(\linewidth - 10\tabcolsep) * \real{0.1111}}
  >{\raggedright\arraybackslash}p{(\linewidth - 10\tabcolsep) * \real{0.1667}}
  >{\raggedright\arraybackslash}p{(\linewidth - 10\tabcolsep) * \real{0.1667}}@{}}
\toprule\noalign{}
\begin{minipage}[b]{\linewidth}\raggedright
\#
\end{minipage} & \begin{minipage}[b]{\linewidth}\raggedright
パラメータ
\end{minipage} & \begin{minipage}[b]{\linewidth}\raggedright
記号
\end{minipage} & \begin{minipage}[b]{\linewidth}\raggedright
値
\end{minipage} & \begin{minipage}[b]{\linewidth}\raggedright
次元
\end{minipage} & \begin{minipage}[b]{\linewidth}\raggedright
変換
\end{minipage} \\
\midrule\noalign{}
\endhead
\bottomrule\noalign{}
\endlastfoot
P5 & 素電荷 & e & 1.602 × 10⁻¹⁹ C（定義値） & \([Q]\) &
電荷の最小単位（量子化の基準） \\
\end{longtable}
}

\begin{itemize}
\tightlist
\item
  \(L, T, M, \Theta\)
  のいずれにも依存しない。\textbf{時空の幾何学とは本質的に独立}。
\item
  電磁気力の結合定数 \(\alpha = e^2/(4\pi\epsilon_0\hbar c)\)
  を構成する。
\item
  TOEでの問い：なぜ電荷は量子化されるのか。ディラックの量子化条件により、磁気単極子が存在すれば電荷の量子化は必然となる。
\end{itemize}

\paragraph{\texorpdfstring{グループ
V：電荷＋時空＋質量（\(M, L, T, Q\)）}{グループ V：電荷＋時空＋質量（M, L, T, Q）}}\label{ux30b0ux30ebux30fcux30d7-vux96fbux8377ux6642ux7a7aux8ceaux91cfm-l-t-q}

{\def\LTcaptype{none} % do not increment counter
\begin{longtable}[]{@{}
  >{\raggedright\arraybackslash}p{(\linewidth - 10\tabcolsep) * \real{0.0833}}
  >{\raggedright\arraybackslash}p{(\linewidth - 10\tabcolsep) * \real{0.3056}}
  >{\raggedright\arraybackslash}p{(\linewidth - 10\tabcolsep) * \real{0.1667}}
  >{\raggedright\arraybackslash}p{(\linewidth - 10\tabcolsep) * \real{0.1111}}
  >{\raggedright\arraybackslash}p{(\linewidth - 10\tabcolsep) * \real{0.1667}}
  >{\raggedright\arraybackslash}p{(\linewidth - 10\tabcolsep) * \real{0.1667}}@{}}
\toprule\noalign{}
\begin{minipage}[b]{\linewidth}\raggedright
\#
\end{minipage} & \begin{minipage}[b]{\linewidth}\raggedright
パラメータ
\end{minipage} & \begin{minipage}[b]{\linewidth}\raggedright
記号
\end{minipage} & \begin{minipage}[b]{\linewidth}\raggedright
値
\end{minipage} & \begin{minipage}[b]{\linewidth}\raggedright
次元
\end{minipage} & \begin{minipage}[b]{\linewidth}\raggedright
変換
\end{minipage} \\
\midrule\noalign{}
\endhead
\bottomrule\noalign{}
\endlastfoot
P6 & 真空の透磁率/誘電率 & μ₀, ε₀ & \(c = 1/\sqrt{\mu_0 \epsilon_0}\) &
\(\mu_0: [M L Q^{-2}]\), \(\epsilon_0: [Q^2 T^2 M^{-1} L^{-3}]\) & 電荷
\(Q\) を時空 \((L,T)\) と質量 \(M\) に接続する \\
\end{longtable}
}

\begin{itemize}
\tightlist
\item
  グループ IV の電荷 \(Q\) を、グループ I〜II の時空・質量と結びつける。
\item
  c
  と従属関係にあり独立ではないが、\textbf{「孤立した電荷が時空とどう接続されるか」}を決める構造的役割を持つ。
\item
  真空が「空」でない（量子的構造を持つ）場合、この定数の意味が変わる可能性がある。
\end{itemize}

\paragraph{次元の階層構造}\label{ux6b21ux5143ux306eux968eux5c64ux69cbux9020}

以上の分類から、基本次元には以下の階層がある。

\[\underbrace{L, T}_{\text{時空（幾何学）}} \xrightarrow{+M} \underbrace{M, L, T}_{\text{物質（量子力学・重力）}} \xrightarrow{+\Theta} \underbrace{M, L, T, \Theta}_{\text{統計（情報）}} \]

\[\underbrace{Q}_{\text{電荷（孤立）}} \xrightarrow{\mu_0, \epsilon_0} \underbrace{M, L, T, Q}_{\text{電磁気学}}\]

自然単位系（\(c = \hbar = k_B = G = 1\)）では、\(L, T, M, \Theta\)
の区別が消去され、すべての物理量が無次元の純粋な数に還元される。ただし電荷
\(Q\) の次元は \(e = 1\) と置くことで別途消去される。

\textbf{このとき残る「数」こそが、万物の理論が説明すべき真の任意パラメータである。}

\textbf{注意:} 2019年のSI単位系改定により、c, ℏ, k\_B, e
はすべて\textbf{定義値}（exact
value）とされている。これは「値が分かった」のではなく、「単位の方をこの定数に合わせた」ということであり、定数の物理的起源が解明されたわけではない。

\subsubsection{1.1
「19個の任意パラメータ」の独立性分析}\label{ux500bux306eux4efbux610fux30d1ux30e9ux30e1ux30fcux30bfux306eux72ecux7acbux6027ux5206ux6790}

以下に、標準模型で「任意パラメータ」と呼ばれている19個 +
重力・宇宙論・ニュートリノのパラメータを列挙し、それぞれが\textbf{真に独立なパラメータか、§1.0
の基本定数や他のパラメータから導出される従属量か}を分析する。

\paragraph{A.
力の結合定数（従来4個とされるもの）}\label{a.-ux529bux306eux7d50ux5408ux5b9aux6570ux5f93ux67654ux500bux3068ux3055ux308cux308bux3082ux306e}

\textbf{\#1 電磁気力の結合定数 α ≈ 1/137.036}

\[\alpha = \frac{e^2}{4\pi\epsilon_0\hbar c}\]

\begin{itemize}
\tightlist
\item
  \textbf{独立か：❌ 従属。} §1.0 の基本定数 P1(c), P2(ℏ), P4(e), P6(ε₀)
  の組み合わせにすぎない。
\item
  α ≈ 1/137 は「神秘的な数」と言われるが、e, ε₀, ℏ, c
  の値を代入すれば自動的にこの値になる。導くも何もない。
\item
  真の問いは「なぜ素電荷 e がこの値か」であり、α は問いの本体ではない。
\item
  \textbf{任意パラメータとしてカウントすべきではない。e（P4）と二重計上。}
\end{itemize}

\textbf{\#2 強い力の結合定数 α\_s ≈ 0.12（91 GeV）}

\[\alpha_s(Q^2) \approx \frac{12\pi}{(33-2N_f)\ln(Q^2/\Lambda_{QCD}^2)} \quad \text{（1ループ近似）}\]

\begin{itemize}
\tightlist
\item
  \textbf{独立か：△ 条件付き。} α\_s の値自体はエネルギースケール Q
  に依存し定数ではない。独立パラメータは Λ\_QCD ≈ 200
  MeV（QCDスケール）1個。
\item
  ただしこの導出式は摂動展開の1ループ近似であり、低エネルギー（大距離）で破綻する。厳密解は未発見（ミレニアム懸賞問題）。
\item
  空間依存性は
  \(1/r^2\)（べき乗則）ではなく対数関数。重力・電磁気力とは質的に異なる。
\item
  \textbf{Λ\_QCD を独立パラメータとして認めるが、α\_s 自体は従属量。}
\end{itemize}

\textbf{\#3 弱い力の結合定数 α\_w ≈ 1/30}

\[\alpha_w = \frac{g^2}{4\pi} , \quad G_F = \frac{\sqrt{2}}{8} \frac{g^2}{M_W^2} = \frac{1}{\sqrt{2}v^2}\]

\begin{itemize}
\tightlist
\item
  \textbf{独立か：△ 条件付き。} フェルミ定数 \(G_F\)
  はヒッグスの真空期待値 v から導出される（\(G_F = 1/(\sqrt{2}v^2)\)）。
\item
  電弱統一理論では、弱い力と電磁気力は同じ起源（SU(2)×U(1)対称性）を持ち、ワインバーグ角
  \(\theta_W\) で混合される。
\item
  \textbf{v（ヒッグスVEV）と \(\theta_W\) が独立パラメータ。α\_w
  自体は従属量。}
\end{itemize}

\textbf{\#4 重力の結合定数 α\_G ≈ 5.9 × 10⁻³⁹}

\[\alpha_G = \frac{Gm_p^2}{\hbar c}\]

\begin{itemize}
\tightlist
\item
  \textbf{独立か：❌ 従属。} §1.0 の P1(c), P2(ℏ), P5(G) と粒子質量
  \(m_p\) の組み合わせ。
\item
  しかも「どの質量 m
  を使うか」が任意（陽子？電子？プランク質量？）。定義すら一意でない。
\item
  \textbf{任意パラメータとしてカウントすべきではない。G（P5）および粒子質量と二重計上。}
\end{itemize}

\textbf{結合定数の小括：}
「4つの結合定数」のうち、独立パラメータとしてカウントできるのは
\textbf{Λ\_QCD} と \textbf{ワインバーグ角 \(\theta_W\)} の2個のみ。α と
α\_G は §1.0 の基本定数と粒子質量の組み合わせであり、独立ではない。

\paragraph{B. 宇宙定数（1個）}\label{b.-ux5b87ux5b99ux5b9aux65701ux500b}

\textbf{\#5 宇宙定数 Λ ≈ 1.1 × 10⁻⁵² m⁻²}

\begin{itemize}
\tightlist
\item
  \textbf{独立か：✅ 独立。}
  他のパラメータから導出する方法は知られていない。
\item
  量子場の零点エネルギーから理論予測すると観測値と120桁乖離する。
\item
  次元は \([L^{-2}]\)（曲率）。時空の幾何学的性質を決める。
\item
  \textbf{真に独立な任意パラメータ。}
\end{itemize}

\paragraph{C.
粒子の質量（9個）}\label{c.-ux7c92ux5b50ux306eux8ceaux91cf9ux500b}

標準模型では、粒子の質量はヒッグス場との湯川結合定数 \(y_i\) を通じて
\(m_i = y_i \cdot v / \sqrt{2}\) で決まる。

{\def\LTcaptype{none} % do not increment counter
\begin{longtable}[]{@{}
  >{\raggedright\arraybackslash}p{(\linewidth - 8\tabcolsep) * \real{0.0484}}
  >{\raggedright\arraybackslash}p{(\linewidth - 8\tabcolsep) * \real{0.1774}}
  >{\raggedright\arraybackslash}p{(\linewidth - 8\tabcolsep) * \real{0.0968}}
  >{\raggedright\arraybackslash}p{(\linewidth - 8\tabcolsep) * \real{0.5645}}
  >{\raggedright\arraybackslash}p{(\linewidth - 8\tabcolsep) * \real{0.1129}}@{}}
\toprule\noalign{}
\begin{minipage}[b]{\linewidth}\raggedright
\#
\end{minipage} & \begin{minipage}[b]{\linewidth}\raggedright
パラメータ
\end{minipage} & \begin{minipage}[b]{\linewidth}\raggedright
質量
\end{minipage} & \begin{minipage}[b]{\linewidth}\raggedright
湯川結合定数 \(y_i = m_i\sqrt{2}/v\)
\end{minipage} & \begin{minipage}[b]{\linewidth}\raggedright
独立か
\end{minipage} \\
\midrule\noalign{}
\endhead
\bottomrule\noalign{}
\endlastfoot
6 & 電子 m\_e & 0.511 MeV & \(y_e \approx 2.9 \times 10^{-6}\) & △
\(y_e\) が独立。m\_e は \(y_e \times v\) \\
7 & ミューオン m\_μ & 105.7 MeV & \(y_\mu \approx 6.1 \times 10^{-4}\) &
同上 \\
8 & タウ m\_τ & 1776.9 MeV & \(y_\tau \approx 1.0 \times 10^{-2}\) &
同上 \\
9 & アップ m\_u & ≈ 2.2 MeV & \(y_u \approx 1.3 \times 10^{-5}\) &
同上 \\
10 & ダウン m\_d & ≈ 4.7 MeV & \(y_d \approx 2.7 \times 10^{-5}\) &
同上 \\
11 & チャーム m\_c & ≈ 1275 MeV & \(y_c \approx 7.3 \times 10^{-3}\) &
同上 \\
12 & ストレンジ m\_s & ≈ 95 MeV & \(y_s \approx 5.5 \times 10^{-4}\) &
同上 \\
13 & トップ m\_t & ≈ 173200 MeV & \(y_t \approx 1.0\)（ほぼ1） & 同上 \\
14 & ボトム m\_b & ≈ 4180 MeV & \(y_b \approx 2.4 \times 10^{-2}\) &
同上 \\
\end{longtable}
}

\begin{itemize}
\tightlist
\item
  \textbf{独立か：△ 条件付き。} 質量 \(m_i\) 自体は
  \(y_i \times v/\sqrt{2}\)
  で導出される従属量。真に独立なのは湯川結合定数
  \(y_i\)（9個）とヒッグスVEV \(v\)（1個）。
\item
  注目すべきは、\(y_t \approx 1\)（トップクォーク）だけが「自然な」値であり、残り8個は
  \(10^{-6}\) 〜 \(10^{-2}\)
  という不自然に小さい値。\textbf{なぜ湯川結合定数にこれほどの階層があるかは未解決。}
\item
  \textbf{9個の質量は、9個の湯川結合定数 \(y_i\) + v
  に分解される。質量として二重計上すべきではない。}
\end{itemize}

\paragraph{D.
CKM行列のパラメータ（4個）}\label{d.-ckmux884cux5217ux306eux30d1ux30e9ux30e1ux30fcux30bf4ux500b}

{\def\LTcaptype{none} % do not increment counter
\begin{longtable}[]{@{}llll@{}}
\toprule\noalign{}
\# & パラメータ & 値 & 独立か \\
\midrule\noalign{}
\endhead
\bottomrule\noalign{}
\endlastfoot
15 & 混合角 θ₁₂ & ≈ 13.0°（カビボ角） & ✅ 独立 \\
16 & 混合角 θ₂₃ & ≈ 2.4° & ✅ 独立 \\
17 & 混合角 θ₁₃ & ≈ 0.2° & ✅ 独立 \\
18 & CP位相 δ\_CKM & ≈ 1.2 rad & ✅ 独立 \\
\end{longtable}
}

\begin{itemize}
\tightlist
\item
  \textbf{独立か：✅ 独立。}
  他のパラメータから導出する方法は知られていない。
\item
  ただし、クォークの世代構造（なぜ3世代か）が解明されれば、幾何学的・群論的に決まる可能性はある。
\item
  \textbf{弱い力に固有のパラメータ。弱い力が存在しなければ不要。}
\end{itemize}

\paragraph{E.
ヒッグス場のパラメータ（2個）}\label{e.-ux30d2ux30c3ux30b0ux30b9ux5834ux306eux30d1ux30e9ux30e1ux30fcux30bf2ux500b}

\textbf{\#19 ヒッグス粒子質量 m\_H ≈ 125.1 GeV}

\[m_H = \sqrt{2\lambda} \cdot v\]

\begin{itemize}
\tightlist
\item
  \textbf{独立か：△ 条件付き。} ヒッグス自己結合定数 λ と VEV v
  から導出される。真に独立なのは λ と v。
\item
  m\_H を「任意パラメータ」とするのは λ × v を一つの数字にまとめただけ。
\end{itemize}

\textbf{\#20 真空期待値 v ≈ 246.2 GeV}

\[v = (\sqrt{2}G_F)^{-1/2} = \frac{2M_W}{g}\]

\begin{itemize}
\tightlist
\item
  \textbf{独立か：✅ 独立。}
  電弱対称性の破れのスケールを決める。全粒子の質量スケールの基準。
\item
  弱い力の到達距離（W/Zボソン質量）もこの値で決まる。
\item
  \textbf{最も重要な独立パラメータの一つ。}
\end{itemize}

\paragraph{F.
QCD真空のパラメータ（1個）}\label{f.-qcdux771fux7a7aux306eux30d1ux30e9ux30e1ux30fcux30bf1ux500b}

\textbf{\#21 QCD真空角 θ\_QCD \textless{} 10⁻¹⁰}

\begin{itemize}
\tightlist
\item
  \textbf{独立か：✅ 独立（ただし値がほぼゼロ）。}
\item
  実験上限 \textless{} 10⁻¹⁰ であり、物理的にはゼロである可能性が高い。
\item
  もしゼロならば何らかの対称性（Peccei-Quinn対称性）が背後にあり、任意パラメータではなく\textbf{構造的にゼロ}。
\item
  \textbf{「任意パラメータ」か「構造的にゼロ」かが未決定。}
\end{itemize}

\paragraph{G.
ニュートリノのパラメータ（7個以上）}\label{g.-ux30cbux30e5ux30fcux30c8ux30eaux30ceux306eux30d1ux30e9ux30e1ux30fcux30bf7ux500bux4ee5ux4e0a}

{\def\LTcaptype{none} % do not increment counter
\begin{longtable}[]{@{}
  >{\raggedright\arraybackslash}p{(\linewidth - 6\tabcolsep) * \real{0.1154}}
  >{\raggedright\arraybackslash}p{(\linewidth - 6\tabcolsep) * \real{0.4231}}
  >{\raggedright\arraybackslash}p{(\linewidth - 6\tabcolsep) * \real{0.1923}}
  >{\raggedright\arraybackslash}p{(\linewidth - 6\tabcolsep) * \real{0.2692}}@{}}
\toprule\noalign{}
\begin{minipage}[b]{\linewidth}\raggedright
\#
\end{minipage} & \begin{minipage}[b]{\linewidth}\raggedright
パラメータ
\end{minipage} & \begin{minipage}[b]{\linewidth}\raggedright
値
\end{minipage} & \begin{minipage}[b]{\linewidth}\raggedright
独立か
\end{minipage} \\
\midrule\noalign{}
\endhead
\bottomrule\noalign{}
\endlastfoot
22 & Δm²₂₁ & ≈ 7.5 × 10⁻⁵ eV² & ✅
独立（質量差の二乗。絶対質量は未測定） \\
23 & Δm²₃₂ & ≈ 2.5 × 10⁻³ eV² & ✅ 独立 \\
24 & PMNS θ₁₂ & ≈ 33.4° & ✅ 独立 \\
25 & PMNS θ₂₃ & ≈ 49° & ✅ 独立 \\
26 & PMNS θ₁₃ & ≈ 8.5° & ✅ 独立 \\
27 & PMNS δ & 未確定 & ✅ 独立（未測定） \\
28 & マヨラナ位相 α₁, α₂ & 未測定 & ？ 存在するかも未確定 \\
\end{longtable}
}

\subsubsection{1.8
独立性分析の総括}\label{ux72ecux7acbux6027ux5206ux6790ux306eux7dcfux62ec}

上記の分析により、「19個の任意パラメータ」の実態は以下の通りである。

\paragraph{従属量（他から導出可能で、独立パラメータとしてカウントすべきでないもの）}\label{ux5f93ux5c5eux91cfux4ed6ux304bux3089ux5c0eux51faux53efux80fdux3067ux72ecux7acbux30d1ux30e9ux30e1ux30fcux30bfux3068ux3057ux3066ux30abux30a6ux30f3ux30c8ux3059ux3079ux304dux3067ux306aux3044ux3082ux306e}

{\def\LTcaptype{none} % do not increment counter
\begin{longtable}[]{@{}
  >{\raggedright\arraybackslash}p{(\linewidth - 4\tabcolsep) * \real{0.3793}}
  >{\raggedright\arraybackslash}p{(\linewidth - 4\tabcolsep) * \real{0.2759}}
  >{\raggedright\arraybackslash}p{(\linewidth - 4\tabcolsep) * \real{0.3448}}@{}}
\toprule\noalign{}
\begin{minipage}[b]{\linewidth}\raggedright
パラメータ
\end{minipage} & \begin{minipage}[b]{\linewidth}\raggedright
導出式
\end{minipage} & \begin{minipage}[b]{\linewidth}\raggedright
二重計上先
\end{minipage} \\
\midrule\noalign{}
\endhead
\bottomrule\noalign{}
\endlastfoot
α ≈ 1/137 & \(e^2/(4\pi\epsilon_0\hbar c)\) & P1(c), P2(ℏ), P4(e),
P6(ε₀) \\
α\_G ≈ 10⁻³⁹ & \(Gm^2/(\hbar c)\) & P1(c), P2(ℏ), P5(G) +
粒子質量。しかも m が任意で定義不一意 \\
α\_s（特定スケールでの値） &
\(12\pi/((33-2N_f)\ln(Q^2/\Lambda_{QCD}^2))\) & Λ\_QCD に帰着 \\
α\_w & \(g^2/(4\pi)\) & v と \(\theta_W\) に帰着 \\
粒子質量 9個 & \(m_i = y_i \cdot v/\sqrt{2}\) & 湯川結合定数 \(y_i\) + v
に帰着 \\
ヒッグス質量 m\_H & \(m_H = \sqrt{2\lambda} \cdot v\) & ヒッグス自己結合
λ + v に帰着 \\
\end{longtable}
}

\paragraph{真に独立なパラメータ}\label{ux771fux306bux72ecux7acbux306aux30d1ux30e9ux30e1ux30fcux30bf}

{\def\LTcaptype{none} % do not increment counter
\begin{longtable}[]{@{}
  >{\raggedright\arraybackslash}p{(\linewidth - 4\tabcolsep) * \real{0.3462}}
  >{\raggedright\arraybackslash}p{(\linewidth - 4\tabcolsep) * \real{0.4231}}
  >{\raggedright\arraybackslash}p{(\linewidth - 4\tabcolsep) * \real{0.2308}}@{}}
\toprule\noalign{}
\begin{minipage}[b]{\linewidth}\raggedright
カテゴリ
\end{minipage} & \begin{minipage}[b]{\linewidth}\raggedright
パラメータ
\end{minipage} & \begin{minipage}[b]{\linewidth}\raggedright
個数
\end{minipage} \\
\midrule\noalign{}
\endhead
\bottomrule\noalign{}
\endlastfoot
基本物理定数 & c, ℏ, k\_B, e, G, (μ₀ε₀) & 6個（ただし μ₀ε₀ は c
に従属） \\
宇宙定数 & Λ & 1個 \\
強い力のスケール & Λ\_QCD & 1個 \\
電弱混合 & ワインバーグ角 \(\theta_W\) & 1個 \\
湯川結合定数 & \(y_e, y_\mu, y_\tau, y_u, y_d, y_c, y_s, y_t, y_b\) &
9個 \\
ヒッグス & 自己結合 λ、真空期待値 v & 2個 \\
CKM行列 & θ₁₂, θ₂₃, θ₁₃, δ\_CKM & 4個 \\
QCD真空 & θ\_QCD & 1個（構造的にゼロの可能性） \\
ニュートリノ & Δm²₂₁, Δm²₃₂, θ₁₂, θ₂₃, θ₁₃, δ, (α₁, α₂) & 7個以上 \\
\textbf{合計} & & \textbf{約32個}（基本定数5個 +
任意パラメータ約27個） \\
\end{longtable}
}

\paragraph{本稿の結論（暫定）}\label{ux672cux7a3fux306eux7d50ux8ad6ux66abux5b9a}

「19個の任意パラメータ」という従来の数え方は、以下の点で不正確である。

\begin{enumerate}
\def\labelenumi{\arabic{enumi}.}
\tightlist
\item
  \textbf{二重計上：} α, α\_G, α\_w, 粒子質量, m\_H
  は他のパラメータの組み合わせであり、独立パラメータとしてカウントすべきでない。
\item
  \textbf{過少計上：} 重力（G,
  Λ）、ニュートリノ（7個以上）、隠れたパラメータ（Λ\_QCD, \(\theta_W\),
  湯川結合定数, λ）を含めていない。
\item
  \textbf{基本定数の除外：} c, ℏ, G, k\_B, e
  を「定数だから」と除外しているが、その値の由来は未解明。
\end{enumerate}

万物の理論のKPIとしては、これら約32個のパラメータについて以下を評価する。

\[\text{KPI} = \frac{\text{理論が内部的に導出した値と実験値が一致したパラメータ数}}{\text{検証対象の総パラメータ数（約32個）}}\]

理想は、最小限の入力（0〜2個）から約32個すべてを再現すること。現状の標準模型のKPIは実質ゼロである（入力をそのまま返しているだけ）。

\begin{center}\rule{0.5\linewidth}{0.5pt}\end{center}

\subsection{2.
結合定数の構造分析}\label{ux7d50ux5408ux5b9aux6570ux306eux69cbux9020ux5206ux6790}

§1.1
で示した通り、「4つの結合定数」はいずれも従属量であるが、その導出式の構造自体に重要な問題が隠されている。

\subsubsection{2.1 電磁気力の結合定数
α}\label{ux96fbux78c1ux6c17ux529bux306eux7d50ux5408ux5b9aux6570-ux3b1}

\[\alpha = \frac{e^2}{4\pi\epsilon_0\hbar c} \approx \frac{1}{137.036}\]

\textbf{構成要素の次元分析：}

{\def\LTcaptype{none} % do not increment counter
\begin{longtable}[]{@{}lll@{}}
\toprule\noalign{}
要素 & 次元 & §1.0 分類 \\
\midrule\noalign{}
\endhead
\bottomrule\noalign{}
\endlastfoot
\(e^2\) & \([Q^2]\) & グループ IV（電荷） \\
\(\epsilon_0\) & \([Q^2 T^2 M^{-1} L^{-3}]\) & グループ
V（電荷↔時空） \\
\(\hbar\) & \([M L^2 T^{-1}]\) & グループ II（時空＋質量） \\
\(c\) & \([L T^{-1}]\) & グループ I（純粋時空） \\
\end{longtable}
}

次元を検証すると：\([Q^2] / ([Q^2 T^2 M^{-1} L^{-3}] \cdot [M L^2 T^{-1}] \cdot [L T^{-1}]) = [1]\)（無次元）✓

\textbf{構造的意味：} α は、電荷（グループ IV）が時空＋質量（グループ I,
II, V）とどう結合するかの比率を表す。5つの基本定数のうち4つ（c, ℏ, e,
ε₀）を使い切って1つの無次元数を作っている。\textbf{α
の値に謎はなく、謎があるとすれば e の値にある。}

\subsubsection{2.2 強い力の結合定数
α\_s}\label{ux5f37ux3044ux529bux306eux7d50ux5408ux5b9aux6570-ux3b1_s}

\[\alpha_s(Q^2) \approx \frac{12\pi}{(33-2N_f)\ln(Q^2/\Lambda_{QCD}^2)} \quad \text{（1ループ近似）}\]

\textbf{構造的問題：}

\begin{enumerate}
\def\labelenumi{\arabic{enumi}.}
\tightlist
\item
  \textbf{定数ではない。} エネルギースケール Q（≒ 空間分解能 ℏ/r
  の逆数）の関数。「結合定数」という名称が誤解を招く。
\item
  \textbf{近似式。}
  1ループ摂動展開の結果であり、厳密解ではない。高次補正を入れるほど複雑化する。
\item
  \textbf{低エネルギーで破綻。} \(Q \to \Lambda_{QCD}\)
  で対数が発散し計算不能。閉じ込め領域は格子QCD（数値計算）に頼る。厳密解の存在はミレニアム懸賞問題として未証明。
\item
  \textbf{空間依存性が幾何学的でない。} 重力・電磁気力の
  \(1/r^2\)（3次元空間での面積の逆数）に対し、α\_s
  は対数関数。次元解析から自然に出る形ではなく、量子補正（仮想粒子対の生成消滅）の結果としてのみ現れる。
\end{enumerate}

\textbf{独立パラメータは Λ\_QCD のみ。} ただし Λ\_QCD
は「カラー荷（離散ラベル）を空間次元 \(L\)
に接続する隠れた次元変換定数」としての役割を持つ（§1.0 で未分類）。

\subsubsection{2.3 弱い力の結合定数
α\_w}\label{ux5f31ux3044ux529bux306eux7d50ux5408ux5b9aux6570-ux3b1_w}

\[\alpha_w = \frac{g^2}{4\pi}, \quad \sin^2\theta_W = 1 - \frac{M_W^2}{M_Z^2} \approx 0.231\]

\[G_F = \frac{1}{\sqrt{2}v^2} \approx 1.166 \times 10^{-5} \text{ GeV}^{-2}\]

\textbf{構造的特徴：}

\begin{enumerate}
\def\labelenumi{\arabic{enumi}.}
\tightlist
\item
  \textbf{電磁気力と統一済み（電弱統一）。} SU(2)×U(1)
  ゲージ理論により、弱い力と電磁気力はワインバーグ角 \(\theta_W\)
  で混合される。低エネルギーでは別の力に見えるが、≈ 100 GeV
  以上で統合される。
\item
  \textbf{到達距離が有限。} W/Zボソンの質量（\(M_W \approx 80.4\) GeV,
  \(M_Z \approx 91.2\) GeV）により、力の到達距離は
  \(\sim \hbar/(M_W c) \approx 10^{-18}\) m
  に制限される。湯川型ポテンシャル \(e^{-M_W r}/r\) に従う。
\item
  \textbf{ヒッグスVEV v に帰着。} フェルミ定数 \(G_F = 1/(\sqrt{2}v^2)\)
  であり、弱い力のスケールはすべて v で決まる。
\end{enumerate}

\textbf{独立パラメータは v と \(\theta_W\)。} α\_w 自体は従属量。

\subsubsection{2.4 重力の結合定数
α\_G}\label{ux91cdux529bux306eux7d50ux5408ux5b9aux6570-ux3b1_g}

\[\alpha_G = \frac{Gm^2}{\hbar c}\]

\textbf{構造的問題：}

\begin{enumerate}
\def\labelenumi{\arabic{enumi}.}
\item
  \textbf{質量 m の選択が任意。}
  他の3つの結合定数は力の性質だけで決まるが、α\_G
  は「どの粒子の質量を使うか」で値が変わる。

  {\def\LTcaptype{none} % do not increment counter
  \begin{longtable}[]{@{}ll@{}}
  \toprule\noalign{}
  質量の選択 & α\_G の値 \\
  \midrule\noalign{}
  \endhead
  \bottomrule\noalign{}
  \endlastfoot
  電子質量 \(m_e\) & ≈ 1.8 × 10⁻⁴⁵ \\
  陽子質量 \(m_p\) & ≈ 5.9 × 10⁻³⁹ \\
  プランク質量 \(m_{Pl}\) & = 1（定義上） \\
  \end{longtable}
  }

  定義すら一意でない「定数」は、定数と呼ぶべきではない。
\item
  \textbf{§1.0 の定数の組み合わせ。} G（P5）, ℏ（P2）,
  c（P1）はすべて基本定数。α\_G
  に新しい情報はなく、粒子質量と基本定数の二重計上にすぎない。
\item
  \textbf{「重力は弱い」は曲率半径の問題。} α\_G
  が小さいことは「重力が弱い」のではなく、時空の曲率半径が電磁気の加速度の曲率半径より圧倒的に大きいことの反映。同じスケールで比較すること自体がカテゴリーエラー。
\end{enumerate}

\subsubsection{2.5
4つの結合定数の比較における根本問題}\label{ux3064ux306eux7d50ux5408ux5b9aux6570ux306eux6bd4ux8f03ux306bux304aux3051ux308bux6839ux672cux554fux984c}

従来の物理学では、4つの結合定数を以下のように並べて「階層性問題」を提起する。

{\def\LTcaptype{none} % do not increment counter
\begin{longtable}[]{@{}lll@{}}
\toprule\noalign{}
力 & 結合定数 & 相対比 \\
\midrule\noalign{}
\endhead
\bottomrule\noalign{}
\endlastfoot
強い力 & α\_s ≈ 1 & 1 \\
電磁気力 & α ≈ 1/137 & 10⁻² \\
弱い力 & α\_w ≈ 1/30 & 10⁻² \\
重力 & α\_G ≈ 10⁻³⁹ & 10⁻³⁹ \\
\end{longtable}
}

しかし、本稿の分析によりこの比較は以下の点で無効である。

\begin{enumerate}
\def\labelenumi{\arabic{enumi}.}
\tightlist
\item
  \textbf{比較対象が同質でない。} α は基本定数の組み合わせ、α\_s
  はスケール依存関数、α\_w は v と θ\_W の従属量、α\_G
  は質量選択に依存する不定量。これらを同じ表に並べること自体に意味がない。
\item
  \textbf{固有スケールが異なる。} 各力は固有の空間スケール（10⁻¹⁸ m 〜
  宇宙スケール）で作用しており、同一のエネルギースケールで比較する根拠がない。
\item
  \textbf{空間依存性が質的に異なる。} 重力・電磁気力は
  \(1/r^2\)（幾何学的）、強い力は対数関数（量子補正）、弱い力は指数減衰（湯川型）。異なる数学的構造を持つ量の「強さ」を単一の数値で比較することは、物理的に無意味。
\item
  \textbf{「階層性問題」は擬似問題。}
  上記1〜3より、「なぜ重力だけが39桁も弱いのか」という問い自体が不適切に構成されている。
\end{enumerate}

\begin{center}\rule{0.5\linewidth}{0.5pt}\end{center}

\subsection{3.
粒子質量の構造分析}\label{ux7c92ux5b50ux8ceaux91cfux306eux69cbux9020ux5206ux6790}

\subsubsection{3.1
湯川結合定数の階層}\label{ux6e6fux5dddux7d50ux5408ux5b9aux6570ux306eux968eux5c64}

§1.1 C で示した通り、粒子質量 \(m_i\) は湯川結合定数 \(y_i\)
とヒッグスVEV \(v\) の積である。

\[m_i = \frac{y_i \cdot v}{\sqrt{2}}\]

真に独立なのは \(y_i\) であり、その値の階層は以下の通り。

{\def\LTcaptype{none} % do not increment counter
\begin{longtable}[]{@{}lllllll@{}}
\toprule\noalign{}
世代 & レプトン & \(y\) & クォーク（上） & \(y\) & クォーク（下） &
\(y\) \\
\midrule\noalign{}
\endhead
\bottomrule\noalign{}
\endlastfoot
第1 & \(e\) & 2.9 × 10⁻⁶ & \(u\) & 1.3 × 10⁻⁵ & \(d\) & 2.7 × 10⁻⁵ \\
第2 & \(\mu\) & 6.1 × 10⁻⁴ & \(c\) & 7.3 × 10⁻³ & \(s\) & 5.5 × 10⁻⁴ \\
第3 & \(\tau\) & 1.0 × 10⁻² & \(t\) & \textbf{≈ 1} & \(b\) & 2.4 ×
10⁻² \\
\end{longtable}
}

\textbf{注目すべき構造：}

\begin{enumerate}
\def\labelenumi{\arabic{enumi}.}
\tightlist
\item
  \textbf{トップクォーク \(y_t \approx 1\) のみが「自然」。} 他の8個は
  \(10^{-6}\) 〜 \(10^{-2}\) と不自然に小さい。なぜ \(y_t\)
  だけがヒッグスVEV のスケールと一致するのかは未解決。
\item
  \textbf{世代間に規則性がある。}
  同じ列（レプトン、上型クォーク、下型クォーク）内で、世代が上がるごとに概ね2〜3桁ずつ増加する。これは偶然か、何らかの構造の反映か。
\item
  \textbf{9個が独立か。}
  世代間の比率に規則性があるなら、独立パラメータは「基準質量 +
  比率規則」に縮約される可能性がある。
\end{enumerate}

\subsubsection{3.2
質量の起源に関する問題}\label{ux8ceaux91cfux306eux8d77ux6e90ux306bux95a2ux3059ux308bux554fux984c}

標準模型では「ヒッグス機構が質量を与える」とされるが、正確にはヒッグス機構は\textbf{質量の仕組み}（対称性の自発的破れ）を提供するだけであり、\textbf{質量の値}（湯川結合定数
\(y_i\)）は何も説明しない。9個の \(y_i\)
がなぜこの値であるかは、標準模型の枠内では原理的に答えられない。

\begin{center}\rule{0.5\linewidth}{0.5pt}\end{center}

\subsection{4.
混合角とCP位相の構造分析}\label{ux6df7ux5408ux89d2ux3068cpux4f4dux76f8ux306eux69cbux9020ux5206ux6790}

\subsubsection{4.1
CKM行列（クォーク）}\label{ckmux884cux5217ux30afux30a9ux30fcux30af}

CKM行列は、弱い力のゲージ固有状態と質量固有状態の間の回転行列である。3世代のユニタリ行列は、3つの回転角と1つの位相で完全に記述される。

{\def\LTcaptype{none} % do not increment counter
\begin{longtable}[]{@{}lll@{}}
\toprule\noalign{}
パラメータ & 値 & 物理的意味 \\
\midrule\noalign{}
\endhead
\bottomrule\noalign{}
\endlastfoot
θ₁₂ ≈ 13.0° & カビボ角 & 第1-2世代間の混合。最大 \\
θ₂₃ ≈ 2.4° & & 第2-3世代間。小さい \\
θ₁₃ ≈ 0.2° & & 第1-3世代間。極めて小さい \\
δ\_CKM ≈ 1.2 rad & CP位相 & 物質・反物質の非対称性 \\
\end{longtable}
}

\textbf{構造的問題：}

\begin{enumerate}
\def\labelenumi{\arabic{enumi}.}
\tightlist
\item
  \textbf{なぜ3世代か。}
  CKM行列が3×3であるのは、クォークが3世代あるからだが、なぜ3世代かは未解明。2世代ならCP位相は不要、4世代以上なら追加のパラメータが必要。
\item
  \textbf{混合角の階層。} θ₁₂ \textgreater{} θ₂₃ \textgreater{} θ₁₃
  という階層がある。これは湯川結合定数の階層（世代間の質量比）と相関している可能性がある。
\item
  \textbf{弱い力に固有。}
  CKM行列は弱い力がなければ不要。弱い力の存在自体が構造要件から導出されるなら、混合角も構造的に決まる可能性がある。
\end{enumerate}

\subsubsection{4.2
PMNS行列（ニュートリノ）}\label{pmnsux884cux5217ux30cbux30e5ux30fcux30c8ux30eaux30ce}

CKM行列のニュートリノ版。構造は同じだが値が大きく異なる。

{\def\LTcaptype{none} % do not increment counter
\begin{longtable}[]{@{}lll@{}}
\toprule\noalign{}
パラメータ & CKM（クォーク） & PMNS（ニュートリノ） \\
\midrule\noalign{}
\endhead
\bottomrule\noalign{}
\endlastfoot
θ₁₂ & 13.0°（小さい） & 33.4°（大きい） \\
θ₂₃ & 2.4°（小さい） & 49°（ほぼ最大混合） \\
θ₁₃ & 0.2°（極小） & 8.5°（小さいが有限） \\
\end{longtable}
}

\textbf{CKMとPMNSの際立った対比：}
クォークの混合は小さく（対角成分が支配的）、ニュートリノの混合は大きい（非対角成分が顕著）。同じ数学的構造（3×3ユニタリ行列）でありながら、なぜ値がこれほど異なるかは未解決。

\begin{center}\rule{0.5\linewidth}{0.5pt}\end{center}

\subsection{5.
ヒッグス場と宇宙定数の構造分析}\label{ux30d2ux30c3ux30b0ux30b9ux5834ux3068ux5b87ux5b99ux5b9aux6570ux306eux69cbux9020ux5206ux6790}

\subsubsection{5.1 ヒッグス場}\label{ux30d2ux30c3ux30b0ux30b9ux5834}

\textbf{独立パラメータ：} 自己結合定数 λ ≈ 0.13、真空期待値 v ≈ 246.2
GeV

\[V(\phi) = -\mu^2|\phi|^2 + \lambda|\phi|^4, \quad v = \frac{\mu}{\sqrt{\lambda}}, \quad m_H = \sqrt{2\lambda}\cdot v\]

\textbf{構造的位置づけ：}

\begin{itemize}
\tightlist
\item
  v
  は電弱対称性の破れのスケールを決め、全粒子質量の基準となる最重要パラメータ。
\item
  λ はヒッグス粒子自身の質量を決めるが、値 ≈ 0.13 に理論的根拠はない。
\item
  v と λ
  の起源は標準模型の枠内では説明できない。「なぜ対称性が破れるか」は構造的に仮定されているだけ。
\end{itemize}

\subsubsection{5.2 宇宙定数 Λ}\label{ux5b87ux5b99ux5b9aux6570-ux3bb}

\[G_{\mu\nu} + \Lambda g_{\mu\nu} = \frac{8\pi G}{c^4} T_{\mu\nu}\]

\textbf{独立パラメータ：} Λ ≈ 1.1 × 10⁻⁵² m⁻²

\paragraph{実測値の導出（幾何学的方法）}\label{ux5b9fux6e2cux5024ux306eux5c0eux51faux5e7eux4f55ux5b66ux7684ux65b9ux6cd5}

アインシュタイン方程式において、宇宙の加速膨張の観測データ（Ia型超新星、CMB、BAO）から得られるエネルギー密度パラメータ
\(\Omega_\Lambda \approx 0.685\) とハッブル定数
\(H_0 \approx 67.4 \; \text{km/s/Mpc}\) を用いると、

\[\Lambda = 3 \Omega_\Lambda \frac{H_0^2}{c^2} \approx 3 \times 0.685 \times \frac{(2.18 \times 10^{-18} \; \text{s}^{-1})^2}{(3.0 \times 10^8 \; \text{m/s})^2} \approx 1.1 \times 10^{-52} \; \text{m}^{-2}\]

ただし、この導出は\textbf{循環論法}である。\(H_0\) と \(\Omega_\Lambda\)
は宇宙の観測曲率から求めた値であり、それを用いて Λ
を求めているに過ぎない。幾何学からΛの値を「導出」しているのではなく、\textbf{観測された曲率にΛという名前を付けているだけ}である。

\paragraph{量子場理論の理論予測値}\label{ux91cfux5b50ux5834ux7406ux8ad6ux306eux7406ux8ad6ux4e88ux6e2cux5024}

量子場理論では、真空は全場のモードの零点エネルギーの総和として有限のエネルギー密度を持つ。各モードの零点エネルギー
\(\frac{1}{2}\hbar\omega\) を、ある紫外カットオフ \(k_{\max}\)
まで積分すると、真空エネルギー密度は

\[\rho_{\text{vac}} = \frac{1}{(2\pi)^3} \int_0^{k_{\max}} \frac{1}{2} \hbar \omega_k \cdot 4\pi k^2 \, dk = \frac{\hbar}{4\pi^2 c} \int_0^{k_{\max}} \frac{1}{2} \sqrt{k^2 c^2 + m^2 c^4 / \hbar^2} \cdot k^2 \, dk\]

質量ゼロの場（\(m = 0\)）かつカットオフをプランクスケール
\(k_{\max} = k_P = \sqrt{\frac{c^3}{\hbar G}}\) とすると、

\[\rho_{\text{vac}} \sim \frac{\hbar c}{16\pi^2} k_P^4 = \frac{c^7}{16\pi^2 \hbar G^2}\]

これを宇宙定数に換算する（\(\Lambda_{\text{QFT}} = \frac{8\pi G}{c^4} \rho_{\text{vac}}\)）と、

\[\Lambda_{\text{QFT}} \sim \frac{c^3}{2\pi \hbar G} = \frac{1}{2\pi \ell_P^2}\]

ここで \(\ell_P = \sqrt{\hbar G / c^3} \approx 1.616 \times 10^{-35}\) m
はプランク長。したがって、

\[\Lambda_{\text{QFT}} \sim \frac{1}{2\pi \times (1.616 \times 10^{-35})^2} \approx 6.1 \times 10^{68} \; \text{m}^{-2}\]

\paragraph{乖離}\label{ux4e56ux96e2}

{\def\LTcaptype{none} % do not increment counter
\begin{longtable}[]{@{}
  >{\raggedright\arraybackslash}p{(\linewidth - 4\tabcolsep) * \real{0.1304}}
  >{\raggedright\arraybackslash}p{(\linewidth - 4\tabcolsep) * \real{0.4783}}
  >{\raggedright\arraybackslash}p{(\linewidth - 4\tabcolsep) * \real{0.3913}}@{}}
\toprule\noalign{}
\begin{minipage}[b]{\linewidth}\raggedright
\end{minipage} & \begin{minipage}[b]{\linewidth}\raggedright
値 (m⁻²)
\end{minipage} & \begin{minipage}[b]{\linewidth}\raggedright
導出方法
\end{minipage} \\
\midrule\noalign{}
\endhead
\bottomrule\noalign{}
\endlastfoot
実測値 & \(1.1 \times 10^{-52}\) &
宇宙の加速膨張の観測データ（循環論法） \\
QFT理論値 & \(\sim 6 \times 10^{68}\) &
真空零点エネルギーのプランクスケール積分 \\
\textbf{乖離} & \textbf{約120桁} & \\
\end{longtable}
}

\textbf{構造的問題：}

\begin{enumerate}
\def\labelenumi{\arabic{enumi}.}
\tightlist
\item
  \textbf{120桁問題。}
  上記の通り、量子場理論の真空エネルギーから計算した理論予測値と観測値が約120桁乖離する。全物理学で最悪の理論予測。この乖離の本質は、量子場理論が「すべてのモードがプランクスケールまで零点エネルギーを持つ」と仮定していることにある。この仮定自体が正しいかどうかは検証されていない。
\item
  \textbf{次元が曲率。} \([\Lambda] = [L^{-2}]\)
  であり、時空の固有曲率を表す。G（\(M\) ↔
  曲率の変換）と同じ幾何学的カテゴリ。
\item
  \textbf{なぜゼロでないか。} 多くの物理学者はΛ =
  0が「自然」と考えていたが、1998年に宇宙の加速膨張が発見され、Λ
  \textgreater{} 0 が確認された。
\item
  \textbf{カットオフ依存性。} QFT理論値はカットオフ \(k_{\max}\)
  の4乗に比例する。プランクスケールの選択に物理的根拠はなく、\(k_{\max}\)
  を変えれば理論値はいくらでも変わる。つまり、120桁の乖離は「理論の予測が外れた」のではなく、\textbf{定式化自体が未完成}であることを示している。
\end{enumerate}

\subsubsection{5.3 QCD真空角
θ\_QCD}\label{qcdux771fux7a7aux89d2-ux3b8_qcd}

\textbf{独立パラメータ：} θ\_QCD \textless{} 10⁻¹⁰

実験的にほぼゼロだが、理論的にはゼロである必要がない。なぜゼロ近傍なのかは「強いCP問題」として未解決。Peccei-Quinn対称性が存在すれば構造的にゼロとなり、その場合はアクシオン粒子が予言される（未発見）。

\begin{center}\rule{0.5\linewidth}{0.5pt}\end{center}

\subsection{6. 分類結果}\label{ux5206ux985eux7d50ux679c}

\subsubsection{6.1
4つの力の定式化水準の比較}\label{ux3064ux306eux529bux306eux5b9aux5f0fux5316ux6c34ux6e96ux306eux6bd4ux8f03}

本稿の分析により、4つの力は理論としての完成度が大きく異なることが明らかになった。

{\def\LTcaptype{none} % do not increment counter
\begin{longtable}[]{@{}
  >{\raggedright\arraybackslash}p{(\linewidth - 8\tabcolsep) * \real{0.1714}}
  >{\raggedright\arraybackslash}p{(\linewidth - 8\tabcolsep) * \real{0.1714}}
  >{\raggedright\arraybackslash}p{(\linewidth - 8\tabcolsep) * \real{0.2571}}
  >{\raggedright\arraybackslash}p{(\linewidth - 8\tabcolsep) * \real{0.2000}}
  >{\raggedright\arraybackslash}p{(\linewidth - 8\tabcolsep) * \real{0.2000}}@{}}
\toprule\noalign{}
\begin{minipage}[b]{\linewidth}\raggedright
項目
\end{minipage} & \begin{minipage}[b]{\linewidth}\raggedright
重力
\end{minipage} & \begin{minipage}[b]{\linewidth}\raggedright
電磁気力
\end{minipage} & \begin{minipage}[b]{\linewidth}\raggedright
弱い力
\end{minipage} & \begin{minipage}[b]{\linewidth}\raggedright
強い力
\end{minipage} \\
\midrule\noalign{}
\endhead
\bottomrule\noalign{}
\endlastfoot
空間依存性 & \(1/r^2\)（厳密） & \(1/r^2\)（厳密） &
\(e^{-Mr}/r\)（湯川型） & 対数（1ループ近似） \\
全距離で計算可能か & ✅ & ✅ & ✅ & ❌（低エネルギーで破綻） \\
厳密解の存在 & ✅ & ✅ & ✅ & ❌（ミレニアム懸賞問題） \\
荷の次元 & \(M\)（連続） & \(Q\)（連続、量子化） & 離散（±1/2） &
離散（3色） \\
荷の離散化の説明 & 未着手 & 未解決（ディラック条件） & 仮定（SU(2)） &
仮定（SU(3)） \\
次元変換定数 & G（明示的） & e, μ₀ε₀（明示的） & v（別枠扱い） &
Λ\_QCD（隠れている） \\
独立パラメータ数 & G, Λ（2個） & e（1個） & v, θ\_W（2個） &
Λ\_QCD（1個） \\
\end{longtable}
}

\subsubsection{6.2
パラメータの依存関係}\label{ux30d1ux30e9ux30e1ux30fcux30bfux306eux4f9dux5b58ux95a2ux4fc2}

以下に、約32個のパラメータの依存関係を示す。矢印は「AからBが導出される」を意味する。

\begin{verbatim}
§1.0 基本定数（独立）
├── c (P1) ─────────────────┐
├── ℏ (P2) ──────────────┐  │
├── k_B (P3) ────────    │  │
├── e (P4) ──────────┐   │  │
├── G (P5) ──────┐   │   │  │
└── μ₀ε₀ (P6) ──┤   │   │  │
                 │   │   │  │
             ┌───┘   │   │  │
             ▼       ▼   ▼  ▼
          α_G ←── m_i   α = e²/(4πε₀ℏc)  ← 従属量（独立でない）
                  ▲
                  │
          y_i × v/√2 = m_i
          ▲       ▲
          │       │
    湯川結合定数  ヒッグスVEV
    y_i (9個)    v (1個)  ← 独立パラメータ
                  │
                  ├── α_w = f(v, θ_W)  ← 従属量
                  ├── G_F = 1/(√2 v²)  ← 従属量
                  └── m_H = √(2λ)·v    ← 従属量
                       ▲
                       │
                  λ (1個)  ← 独立パラメータ

その他の独立パラメータ：
├── Λ_QCD (1個) ── → α_s(Q²) = f(Λ_QCD, Q)  ← 従属量
├── θ_W (1個)
├── Λ (1個) ── 宇宙定数
├── θ_QCD (1個)
├── CKM: θ₁₂, θ₂₃, θ₁₃, δ (4個)
└── PMNS: Δm²₂₁, Δm²₃₂, θ₁₂, θ₂₃, θ₁₃, δ, (α₁, α₂) (7個+)
\end{verbatim}

\subsubsection{6.3
TOEのKPI（受け入れ試験仕様）}\label{toeux306ekpiux53d7ux3051ux5165ux308cux8a66ux9a13ux4ed5ux69d8}

万物の理論の評価基準を以下のように定義する。

\textbf{テストケース：} 約32個の独立パラメータ

{\def\LTcaptype{none} % do not increment counter
\begin{longtable}[]{@{}
  >{\raggedright\arraybackslash}p{(\linewidth - 6\tabcolsep) * \real{0.1333}}
  >{\raggedright\arraybackslash}p{(\linewidth - 6\tabcolsep) * \real{0.3000}}
  >{\raggedright\arraybackslash}p{(\linewidth - 6\tabcolsep) * \real{0.3667}}
  >{\raggedright\arraybackslash}p{(\linewidth - 6\tabcolsep) * \real{0.2000}}@{}}
\toprule\noalign{}
\begin{minipage}[b]{\linewidth}\raggedright
ID
\end{minipage} & \begin{minipage}[b]{\linewidth}\raggedright
カテゴリ
\end{minipage} & \begin{minipage}[b]{\linewidth}\raggedright
パラメータ
\end{minipage} & \begin{minipage}[b]{\linewidth}\raggedright
個数
\end{minipage} \\
\midrule\noalign{}
\endhead
\bottomrule\noalign{}
\endlastfoot
T01〜T05 & 基本定数 & c, ℏ, k\_B, e, G & 5 \\
T06 & 宇宙定数 & Λ & 1 \\
T07 & QCDスケール & Λ\_QCD & 1 \\
T08 & 電弱混合 & θ\_W & 1 \\
T09〜T17 & 湯川結合定数 &
\(y_e, y_\mu, y_\tau, y_u, y_d, y_c, y_s, y_t, y_b\) & 9 \\
T18〜T19 & ヒッグス & λ, v & 2 \\
T20〜T23 & CKM行列 & θ₁₂, θ₂₃, θ₁₃, δ\_CKM & 4 \\
T24 & QCD真空 & θ\_QCD & 1 \\
T25〜T31 & ニュートリノ & Δm²₂₁, Δm²₃₂, θ₁₂, θ₂₃, θ₁₃, δ\_PMNS, (α₁, α₂)
& 7+ \\
& \textbf{合計} & & \textbf{約32} \\
\end{longtable}
}

\textbf{評価方法：}

\[\text{KPI} = \frac{\text{理論が導出した値と実験値が有効桁数内で一致した数}}{32}\]

\textbf{評価基準：}

{\def\LTcaptype{none} % do not increment counter
\begin{longtable}[]{@{}ll@{}}
\toprule\noalign{}
スコア & 評価 \\
\midrule\noalign{}
\endhead
\bottomrule\noalign{}
\endlastfoot
0/32 & 現状の標準模型（入力をそのまま返すだけ） \\
1〜5/32 & 部分的成功（例：θ\_QCD = 0 の説明、結合定数の統一） \\
6〜15/32 & 重要な進展（例：湯川結合定数の階層の説明） \\
16〜25/32 & 大統一理論レベル \\
26〜31/32 & 万物の理論の候補 \\
\textbf{32/32} & \textbf{万物の理論の完成} \\
\end{longtable}
}

\textbf{入力パラメータ数による加点：}

{\def\LTcaptype{none} % do not increment counter
\begin{longtable}[]{@{}ll@{}}
\toprule\noalign{}
入力数 & 評価 \\
\midrule\noalign{}
\endhead
\bottomrule\noalign{}
\endlastfoot
32個（全部入力） & 何も導出していない（現状） \\
10〜31個 & 部分的縮約 \\
3〜9個 & 大幅な縮約 \\
1〜2個 & 優れたTOE候補 \\
\textbf{0個} & \textbf{究極のTOE（初期条件すら不要）} \\
\end{longtable}
}

\begin{center}\rule{0.5\linewidth}{0.5pt}\end{center}

\subsection{7. 結論}\label{ux7d50ux8ad6}

\subsubsection{7.1
「19個の任意パラメータ」の再評価}\label{ux500bux306eux4efbux610fux30d1ux30e9ux30e1ux30fcux30bfux306eux518dux8a55ux4fa1}

本稿の分析により、以下が明らかになった。

\begin{enumerate}
\def\labelenumi{\arabic{enumi}.}
\item
  \textbf{従来の「19個」は不正確である。} 二重計上（α, α\_G, α\_w,
  粒子質量, m\_H）と過少計上（G, Λ, Λ\_QCD, θ\_W, 湯川結合定数, λ,
  ニュートリノ）の両方が存在する。正確には約32個の独立パラメータがある。
\item
  \textbf{基本物理定数の除外は不当である。} c, ℏ, k\_B, e, G
  を「定数だから」と除外しているが、これらの値の物理的起源は未解明であり、TOEの検証対象に含めるべきである。
\item
  \textbf{4つの力は対等に定式化されていない。}
  重力と電磁気力は厳密な解析解（\(1/r^2\)）と明示的な次元変換定数を持つが、強い力は近似解（対数関数）しか持たず低エネルギーで破綻し、弱い力の次元変換定数はヒッグス場に隠されている。
\item
  \textbf{α ≈ 1/137 に謎はない。} \(e^2/(4\pi\epsilon_0\hbar c)\)
  に値を代入すれば自動的に得られる従属量であり、「導けたらすごい」という評価は問いの立て方が間違っている。問うべきは素電荷
  e の値の起源である。
\item
  \textbf{「階層性問題」は擬似問題の可能性がある。}
  4つの結合定数は、定義も、空間依存性も、荷の次元も異なる。同一スケールで比較して「重力だけが弱い」と嘆くのは、異なるカテゴリの量を無理に同一基準で比較した結果である。
\end{enumerate}

\subsubsection{7.2 万物の理論のKPI ---
検証可能性による優先順位}\label{ux4e07ux7269ux306eux7406ux8ad6ux306ekpi-ux691cux8a3cux53efux80fdux6027ux306bux3088ux308bux512aux5148ux9806ux4f4d}

約32個の独立パラメータすべてを一度に検証することは、現段階では不可能である。理由は、4つの力の理論的成熟度が根本的に異なるためである（§6.1参照）。

\paragraph{検証可能な領域：重力場と電磁場の統合}\label{ux691cux8a3cux53efux80fdux306aux9818ux57dfux91cdux529bux5834ux3068ux96fbux78c1ux5834ux306eux7d71ux5408}

重力と電磁気力は、以下の点で検証に必要な条件を満たしている。

\begin{itemize}
\tightlist
\item
  両方とも厳密な解析解（\(1/r^2\)）を持つ
\item
  荷の次元が明示的に定義されている（質量 \(M\)、電荷 \(Q\)）
\item
  次元変換定数が明示的である（\(G\), \(e\), \(\mu_0\epsilon_0\)）
\item
  関与する独立パラメータが特定されている（\(c\), \(\hbar\), \(G\),
  \(e\), \(\mu_0\epsilon_0\), \(\Lambda\)）
\end{itemize}

したがって、統合理論が導出すべきパラメータ間の関係が明確であり、実験値との照合による合否判定が可能である。

\textbf{重力＋電磁気力に関するKPI（検証可能）：}

{\def\LTcaptype{none} % do not increment counter
\begin{longtable}[]{@{}
  >{\raggedright\arraybackslash}p{(\linewidth - 4\tabcolsep) * \real{0.1818}}
  >{\raggedright\arraybackslash}p{(\linewidth - 4\tabcolsep) * \real{0.4091}}
  >{\raggedright\arraybackslash}p{(\linewidth - 4\tabcolsep) * \real{0.4091}}@{}}
\toprule\noalign{}
\begin{minipage}[b]{\linewidth}\raggedright
ID
\end{minipage} & \begin{minipage}[b]{\linewidth}\raggedright
検証項目
\end{minipage} & \begin{minipage}[b]{\linewidth}\raggedright
合格基準
\end{minipage} \\
\midrule\noalign{}
\endhead
\bottomrule\noalign{}
\endlastfoot
T01 & \(G\) と \(e\) の関係 & 一方から他方を導出し、実験値と一致 \\
T02 & \(\alpha = e^2/(4\pi\epsilon_0\hbar c)\) の値 & 構造から \(e\)
を導出し、\(\alpha \approx 1/137.036\) と一致 \\
T03 & \(\Lambda\) の値 & 理論から
\(\Lambda \approx 1.1 \times 10^{-52}\) m⁻² を導出 \\
T04 & 質量 \(M\) と電荷 \(Q\) の関係 & 電荷の量子化条件を導出 \\
\end{longtable}
}

\paragraph{現段階で検証困難な領域：弱い力と強い力}\label{ux73feux6bb5ux968eux3067ux691cux8a3cux56f0ux96e3ux306aux9818ux57dfux5f31ux3044ux529bux3068ux5f37ux3044ux529b}

弱い力と強い力については、「統合できた」ことの検証基準自体が不明確である。

\begin{itemize}
\tightlist
\item
  \textbf{強い力：}
  厳密解が存在しない（ミレニアム懸賞問題）。低エネルギーで摂動論が破綻し、結合定数
  \(\alpha_s\)
  の空間依存性は1ループ近似でしか記述されていない。統合理論の出力と比較すべき「正解」が、近似値しかない。
\item
  \textbf{弱い力：}
  荷（弱アイソスピン）は離散ラベルであり、SI次元を持たない。次元変換定数はヒッグス場の真空期待値
  \(v\) に隠されており、独立パラメータの切り分け自体が曖昧である。
\end{itemize}

未完成の理論を「統合した」と主張しても、それは\textbf{未定式化のものを未定式化のまま束ねただけ}であり、検証のしようがない。

\textbf{結論：}
万物の理論の最初の検証ステップは、重力場と電磁場の統合理論の構築と、上記T01〜T04の検証である。弱い力・強い力の統合は、それぞれの理論自体の厳密な定式化が完了した後に初めて検証対象となる。

\subsubsection{7.3 本稿の限界}\label{ux672cux7a3fux306eux9650ux754c}

本稿は既存のパラメータの分類と問題の整理を行ったものであり、パラメータの導出そのものは行っていない。各パラメータの構造的導出は、前稿
{[}1{]}
で定式化した構造要件を具体的な幾何学に実装する段階で初めて可能となる。

\begin{center}\rule{0.5\linewidth}{0.5pt}\end{center}

\subsection{参考文献}\label{ux53c2ux8003ux6587ux732e}

{[}1{]} 木原範昭, 「万物の理論が満たすべき構造要件の定式化」, 2026. DOI:
\href{https://doi.org/10.5281/zenodo.19601592}{10.5281/zenodo.19601592}

{[}2{]} Particle Data Group, ``Review of Particle Physics'', \emph{Phys.
Rev.~D} \textbf{110}, 030001 (2024).

{[}3{]} S. Weinberg, ``A Model of Leptons'', \emph{Phys. Rev.~Lett.}
\textbf{19}, 1264 (1967).

{[}4{]} P. W. Higgs, ``Broken Symmetries and the Masses of Gauge
Bosons'', \emph{Phys. Rev.~Lett.} \textbf{13}, 508 (1964).

{[}5{]} G. 't Hooft, ``Naturalness, Chiral Symmetry, and Spontaneous
Chiral Symmetry Breaking'', \emph{NATO Adv. Study Inst. Ser. B Phys.}
\textbf{59}, 135 (1980).

{[}6{]} A. Einstein, ``Die Feldgleichungen der Gravitation'',
\emph{Sitzungsberichte der Preussischen Akademie der Wissenschaften zu
Berlin}, 844-847 (1915).

{[}7{]} D. J. Gross and F. Wilczek, ``Ultraviolet Behavior of
Non-Abelian Gauge Theories'', \emph{Phys. Rev.~Lett.} \textbf{30}, 1343
(1973).

{[}8{]} Planck Collaboration, ``Planck 2018 results. VI. Cosmological
parameters'', \emph{Astron. Astrophys.} \textbf{641}, A6 (2020).

\end{document}
